{\small\textbf{The twist (Diego, 1:15) [16:25–17:40]}}\par\smallskip
Here is the twist: the storm made it the \textbf{cheapest} moment.
\begin{itemize}\setlength\itemsep{0pt}
\item \textcolor{navy}{\textbf{[def]}} Look at headroom, the capital banks hold above their requirements. At the end of 2022 the four large banks had about 36 billion euros of it.
\item The full two-percent buffer is 6.7 billion, so it absorbs at most about \textbf{a fifth} of that headroom, and that is an upper bound.
\end{itemize}
\par\smallskip
Bank by bank, the left panel shows the same: ING and Rabobank each had well over ten billion euros of headroom, ABN AMRO about seven. Headroom did shrink between 2021 and 2022, but because risk-weighted assets grew and banks paid out capital, not because of the buffer, which did not yet bind.
\par\smallskip
Why so cheap? Because rising rates lifted net interest income. Profitable banks can build buffers from retained earnings instead of cutting loans.
\par\smallskip
The ECB reached the same conclusion. Its research finds that a gradual build-up with profitable banks "limits these economic costs", and that activating buffers during the 2022–23 tightening "can mitigate implementation costs".
\par\smallskip
So: not the wrong moment. Arguably the right one.
